
\documentclass{beamer}

% Selecting a presentation theme
\usetheme{Madrid}

% Packages
\usepackage{amsmath}
\usepackage{graphicx}

% Presentation details
\title{Introduction to Computer Science}
\subtitle{A Sample Beamer Presentation}
\author{Student Name}
\institute{MNIT Jaipur}
\date{\today}

\begin{document}

% Title Slide
\begin{frame}
\titlepage
\end{frame}

% Introduction Slide
\begin{frame}{Introduction}
Computer Science is the study of computers and computational systems.

\begin{itemize}
    \item Programming
    \item Data Structures
    \item Algorithms
    \item Computer Networks
\end{itemize}
\end{frame}

% Blocks Slide
\begin{frame}{Important Concepts}

\begin{block}{Definition}
An algorithm is a sequence of steps used to solve a problem.
\end{block}

\begin{alertblock}{Important}
Every algorithm should have a clear objective.
\end{alertblock}

\begin{exampleblock}{Example}
A program to add two numbers.
\end{exampleblock}

\end{frame}

% Mathematical Equation Slide
\begin{frame}{Mathematical Expression}

The quadratic equation is:

\[
x = \frac{-b \pm \sqrt{b^2-4ac}}{2a}
\]

This equation is used to calculate the roots of a quadratic equation.

\end{frame}

% Columns Slide
\begin{frame}{Types of Software}

\begin{columns}

\begin{column}{0.5\textwidth}
\textbf{System Software}
\begin{itemize}
    \item Operating Systems
    \item Device Drivers
\end{itemize}
\end{column}

\begin{column}{0.5\textwidth}
\textbf{Application Software}
\begin{itemize}
    \item Web Browsers
    \item Text Editors
\end{itemize}
\end{column}

\end{columns}

\end{frame}

% Overlay Slide
\begin{frame}{Learning Outcomes}

\begin{itemize}
    \item<1-> Understand Beamer.
    \item<2-> Create professional slides.
    \item<3-> Use themes and blocks.
    \item<4-> Display mathematical equations.
\end{itemize}

\end{frame}

% Conclusion Slide
\begin{frame}{Conclusion}
Beamer makes it easy to create structured and professional presentations using LaTeX.

\centering
\textbf{Thank You!}

\end{frame}

\end{document}
